\documentclass[11pt]{article}
\usepackage[a4paper,margin=2.8cm]{geometry}
\usepackage{amsmath,amssymb,amsthm}
\usepackage{enumitem}
\usepackage{xcolor}
\usepackage[colorlinks=true,linkcolor=blue,citecolor=blue]{hyperref}

\newtheorem{theorem}{Theorem}[section]
\newtheorem{lemma}[theorem]{Lemma}
\newtheorem{corollary}[theorem]{Corollary}
\newtheorem{definition}[theorem]{Definition}
\newtheorem{remark}[theorem]{Remark}

\newcommand{\E}{\mathbb{E}}
\renewcommand{\Pr}{\mathbb{P}}
\newcommand{\N}{\mathbb{N}}
\newcommand{\R}{\mathbb{R}}
\newcommand{\Ical}{\mathcal{I}}
\newcommand{\Fcal}{\mathcal{F}}
\newcommand{\ceil}[1]{\left\lceil #1 \right\rceil}

\title{Rumor spreading by uniform push on the complete graph:\\
a formalization-oriented proof of the $O(\log n)$ bound}
\author{Aakash Kumar, Maria Sofia Bucarelli, Emanuele Natale}
\date{}

\begin{document}
\maketitle

\begin{abstract}
\noindent
We give a complete and elementary proof that the uniform \emph{push} protocol
on the complete graph $K_n$, started from a single informed node, informs all
$n$ nodes within $O(\log n)$ rounds with high probability.  The proof is
deliberately optimized for formalization in a proof assistant (Lean~4 +
Mathlib): the probability space is finite and uniform; the only analytic
ingredients are Bernoulli's inequality and Markov's inequality; concentration
in the growth phase is obtained from a one-line exponential-moment induction
for adaptive Bernoulli trials, avoiding both Chernoff bounds for sums of
independent variables and stochastic-domination couplings.  No use is made of
the exponential function, of continuous probability, or of martingale theory.
Constants are deliberately crude.

\smallskip
\noindent
This argument has been \emph{fully formalized} (no \texttt{sorry}) in the
accompanying Lean development (\texttt{RumorSpread/}); the main theorem is
\texttt{RumorPush.push\_informs\_all\_whp}.
\end{abstract}

\section{The model}\label{sec:model}

Fix $n \ge 2$ and let $V = \{0,1,\dots,n-1\}$ be the vertex set of the
complete graph $K_n$.

\begin{definition}[Round randomness]\label{def:round}
A \emph{round configuration} is a function
$\omega \colon V \to V$ with $\omega(v) \ne v$ for all $v$.
Let $R_n$ denote the (finite) set of round configurations; it carries the
uniform probability measure, i.e.\ every node picks a uniformly random
\emph{target} among the other $n-1$ nodes, independently of all other nodes.
\end{definition}

\begin{definition}[Push step]\label{def:step}
For an \emph{informed set} $I \subseteq V$ and a round configuration
$\omega \in R_n$, the next informed set is
\[
  \mathrm{step}(I,\omega) \;=\; I \,\cup\, \omega(I)
  \;=\; I \cup \{\omega(v) : v \in I\}.
\]
(Only informed nodes actually transmit; letting \emph{all} nodes draw a target
and ignoring the choices of uninformed nodes is an equivalent and
formalization-friendlier convention.)
\end{definition}

\begin{definition}[Process]\label{def:process}
Fix a horizon $T \in \N$ and a start vertex $v_0 \in V$. The sample space is
\[
  \Omega \;=\; R_n^{\,T}
  \quad\text{(uniform measure, i.e.\ i.i.d.\ uniform rounds).}
\]
For $\omega = (\omega_0,\dots,\omega_{T-1}) \in \Omega$ define
$I_0(\omega) = \{v_0\}$ and
$I_{t+1}(\omega) = \mathrm{step}(I_t(\omega), \omega_t)$ for $0 \le t < T$.
We write $m_t = |I_t|$ and $U_t = n - m_t$ for the numbers of informed and
uninformed nodes.  Since $\Omega$ is finite, all probabilities are finite
ratios and all expectations are finite sums; ``conditioning on
$I_t$'' means partitioning $\Omega$ according to the first $t$ coordinates.
\end{definition}

Two trivial but load-bearing facts, both immediate from
Definition~\ref{def:step}:

\begin{lemma}[Monotonicity]\label{lem:mono}
$I_t(\omega) \subseteq I_{t+1}(\omega)$ for all $t$ and all $\omega$; hence
$m_t$ is non-decreasing and $U_t$ is non-increasing in $t$, pointwise on
$\Omega$.
\end{lemma}

\begin{lemma}[At most doubling]\label{lem:atmost}
$|I_{t+1}(\omega) \setminus I_t(\omega)| \le |I_t(\omega)|$, i.e.\ the number
$N_t := m_{t+1}-m_t$ of newly informed nodes in one round satisfies
$0 \le N_t \le m_t$.
\end{lemma}

\section{Main theorem}

Throughout, $\log$ denotes the natural logarithm.

\begin{theorem}[Push protocol informs $K_n$ in $O(\log n)$ rounds w.h.p.]
\label{thm:main}
Let $n \ge 2$, let $v_0 \in V$, and set
\[
  T_1 \;=\; \ceil{117 \log n} + 23, \qquad
  T_2 \;=\; \ceil{6 \log n}, \qquad
  T \;=\; T_1 + T_2 .
\]
Then
\[
  \Pr\bigl[\, I_T = V \,\bigr] \;\ge\; 1 - \frac{2}{n}.
\]
In particular all nodes are informed after $O(\log n)$ rounds with high
probability.
\end{theorem}

\begin{remark}
The constants are deliberately crude: the proof is organized to minimize the
analytic toolkit, not the round count.  The sharp answer is
$\log_2 n + \log n + o(\log n)$ (Frieze--Grimmett 1985, Pittel 1987).
Replacing $2/n$ by $n^{-\alpha}$ for any fixed $\alpha \ge 1$ only changes the
constants in $T_1, T_2$.
\end{remark}

The proof has two phases, glued by Lemma~\ref{lem:mono}:
\begin{itemize}[leftmargin=2em]
\item \textbf{Growth phase} (Section~\ref{sec:growth}): while $m_t \le n/2$,
  each round multiplies $m_t$ by $\ge 9/8$ with probability $\ge 1/8$;
  an exponential-moment bound for adaptive trials shows that after $T_1$
  rounds, $m_{T_1} > n/2$ except with probability $\le 1/n$.
\item \textbf{Saturation phase} (Section~\ref{sec:saturation}): once
  $m_t \ge n/2$, the expected number of uninformed nodes contracts by a factor
  $2/3$ per round; after $T_2$ further rounds Markov's inequality gives
  $U_T = 0$ except with probability $\le 1/n$.
\end{itemize}

\section{Elementary inequalities}\label{sec:ineq}

Everything in this section is deterministic real arithmetic.  We write
$x := \tfrac{1}{n-1}$, so $0 < x \le 1$ for $n \ge 2$.

\begin{lemma}[Bernoulli]\label{lem:bernoulli}
For all real $y \ge -1$ (we only use $y\in[-1,1]$) and $m \in \N$:
$(1+y)^m \ge 1 + my$.
\end{lemma}

\begin{proof}
Induction on $m$; available in Mathlib as \texttt{one\_add\_mul\_le\_pow}.
\end{proof}

\begin{lemma}[Second-order lower bound, ``Bonferroni'']\label{lem:bonferroni}
For $0 \le x \le 1$ and $m \in \N$:
\[
  1-(1-x)^m \;\ge\; mx - \binom{m}{2}x^2 \;\ge\; mx\Bigl(1-\frac{mx}{2}\Bigr).
\]
\end{lemma}

\begin{proof}
Induction on $m$. The case $m=0$ is trivial. For the step, using
Lemma~\ref{lem:bernoulli} in the form $(1-x)^m \ge 1-mx$:
\[
  1-(1-x)^{m+1} = \bigl(1-(1-x)^m\bigr) + x(1-x)^m
  \ge \Bigl(mx - \tbinom{m}{2}x^2\Bigr) + x(1-mx)
  = (m+1)x - \tbinom{m+1}{2}x^2 ,
\]
since $\binom{m}{2}+m = \binom{m+1}{2}$. The second inequality is
$\binom{m}{2} \le m^2/2$.
\end{proof}

\begin{lemma}[Upper bound without $e$]\label{lem:upper}
For $0 \le x \le 1$ and $m \in \N$:
\[
  (1-x)^m \;\le\; \frac{1}{1+mx}.
\]
\end{lemma}

\begin{proof}
$(1-x)(1+x) = 1-x^2 \le 1$, so $(1-x)^m (1+x)^m \le 1$, and
$(1+x)^m \ge 1+mx > 0$ by Lemma~\ref{lem:bernoulli}; divide.
\end{proof}

\begin{lemma}[Logarithm bounds]\label{lem:logbounds}
For $x > 0$: $\;1 - \tfrac1x \le \log x \le x - 1$.
\end{lemma}

\begin{proof}
The right inequality is standard (concavity of $\log$, in Mathlib as
\texttt{Real.log\_le\_sub\_one\_of\_pos}); the left follows by applying it to
$1/x$.  These two bounds, together with Mathlib's numeric bound
$\log 2 < 0.6931471808$, are the only facts about $\log$ used in the numeric
verification of Lemma~\ref{lem:phase1}; in particular
$\log\tfrac98 \ge \tfrac19$, $\log\tfrac{16}{15} \ge \tfrac1{16}$ and
$\log\tfrac32 \ge \tfrac13$.
\end{proof}

\begin{lemma}[Markov's inequality, finite uniform version]\label{lem:markov}
Let $\Omega$ be a finite nonempty set with the uniform measure, $f \colon
\Omega \to \R$ with $f \ge 0$, and $a > 0$.  Then
$\Pr[f \ge a] \le \E[f]/a$.  In particular, if $f$ takes values in $\N$, then
$\Pr[f \ne 0] \le \E[f]$.
\end{lemma}

\begin{lemma}[Reverse Markov for bounded variables]\label{lem:revmarkov}
Let $f \colon \Omega \to \R$ with $0 \le f \le b$ pointwise, $b > 0$, and let
$0 < a < b$.  Then
\[
  \Pr[f \ge a] \;\ge\; \frac{\E[f] - a}{b}.
\]
\end{lemma}

\begin{proof}
$\E[f] \le a\,\Pr[f<a] + b\,\Pr[f\ge a] \le a + b\,\Pr[f \ge a]$.
\end{proof}

\section{One-round estimates}\label{sec:oneround}

Fix an informed set $I$ with $|I| = m \ge 1$ and let $\omega \in R_n$ be a
uniform round configuration.  All statements in this section concern this
single round (one uniform draw from $R_n$); by independence of the rounds they
apply verbatim to the conditional law of round $t$ given $I_t = I$.

\begin{lemma}[Contact probability]\label{lem:contact}
For any fixed $u \notin I$,
\[
  \Pr_\omega\bigl[u \in \mathrm{step}(I,\omega)\bigr]
  \;=\; 1 - \Bigl(1-\frac{1}{n-1}\Bigr)^{m}.
\]
\end{lemma}

\begin{proof}
$u \notin \mathrm{step}(I,\omega)$ iff $\omega(v) \ne u$ for every $v \in I$.
The coordinates $(\omega(v))_{v \in I}$ are independent and uniform on
$V \setminus \{v\}$, and $u \ne v$ for $v\in I$, so each misses $u$ with
probability exactly $1 - \frac1{n-1}$; multiply over $v \in I$.
(Formally: count the functions $\omega$ avoiding $u$ on $I$; the count
factorizes coordinate-by-coordinate.)
\end{proof}

\begin{lemma}[Expected growth while below half]\label{lem:expgrowth}
Let $N(\omega) = |\mathrm{step}(I,\omega)| - m$ be the number of newly
informed nodes.  If $1 \le m \le n/2$, then
\[
  \E_\omega[N] \;\ge\; \frac{m}{4}.
\]
\end{lemma}

\begin{proof}
Write $x = \frac1{n-1}$. By linearity of expectation over the $n-m$
uninformed nodes and Lemma~\ref{lem:contact},
$\E[N] = (n-m)\bigl(1-(1-x)^m\bigr)$.
Since $m \le n/2$ and $n \ge 2$ we have
$mx = \frac{m}{n-1} \le \frac{n}{2(n-1)} \le 1$,
so Lemma~\ref{lem:bonferroni} gives
$1-(1-x)^m \ge mx(1-\tfrac{mx}{2}) \ge \tfrac{mx}{2}$.
Using $n - m \ge n/2$:
\[
  \E[N] \;\ge\; \frac{n}{2}\cdot\frac{m}{2(n-1)}
  \;=\; \frac{m}{4}\cdot\frac{n}{n-1} \;\ge\; \frac{m}{4}. \qedhere
\]
\end{proof}

\begin{lemma}[A good round has constant probability]\label{lem:goodround}
If $1 \le m \le n/2$, then
\[
  \Pr_\omega\Bigl[\,|\mathrm{step}(I,\omega)| \ge \tfrac{9}{8}\,m\,\Bigr]
  \;\ge\; \frac18 .
\]
\end{lemma}

\begin{proof}
$N$ satisfies $0 \le N \le m$ (Lemma~\ref{lem:atmost}) and $\E[N] \ge m/4$
(Lemma~\ref{lem:expgrowth}).  By Lemma~\ref{lem:revmarkov} with $a = m/8$,
$b = m$:
\[
  \Pr\bigl[N \ge \tfrac{m}{8}\bigr] \;\ge\; \frac{m/4 - m/8}{m} = \frac18 .
  \qedhere
\]
\end{proof}

\begin{lemma}[Contraction above half]\label{lem:contract}
If $m \ge n/2$, then for every fixed $u \notin I$,
\[
  \Pr_\omega\bigl[u \notin \mathrm{step}(I,\omega)\bigr]
  \;=\; \Bigl(1-\frac1{n-1}\Bigr)^m \;\le\; \frac23 ,
\]
and consequently, with $U' := n - |\mathrm{step}(I,\omega)|$,
\[
  \E_\omega[U'] \;\le\; \frac23\,(n-m).
\]
\end{lemma}

\begin{proof}
By Lemma~\ref{lem:upper}, $(1-x)^m \le \frac{1}{1+mx}$, and
$mx = \frac{m}{n-1} \ge \frac{n}{2(n-1)} \ge \frac12$, so
$(1-x)^m \le \frac{1}{3/2} = \frac23$.  Summing over the $n-m$ uninformed
nodes gives the expectation bound.
\end{proof}

\section{Adaptive Bernoulli trials: an exponential-moment bound}
\label{sec:adaptive}

The only ``concentration'' tool we need.  We state it for our concrete finite
product space; $\Fcal_t$-measurable simply means ``a function of the first $t$
coordinates of $\omega \in \Omega = R_n^T$''.

\begin{lemma}[Exponential moment for adaptive trials]\label{lem:adaptive}
Let $X_0, \dots, X_{T-1} \colon \Omega \to \{0,1\}$ be such that $X_t$ is a
function of $(\omega_0,\dots,\omega_t)$, and suppose there is $p \in [0,1]$
such that for every $t$ and every fixed prefix
$(\omega_0,\dots,\omega_{t-1})$,
\[
  \Pr\bigl[X_t = 1 \,\big|\, \omega_0,\dots,\omega_{t-1}\bigr] \;\ge\; p .
\]
Let $G = \sum_{t<T} X_t$.  Then for every $s \in (0,1]$,
\[
  \E\bigl[s^{\,G}\bigr] \;\le\; \bigl(1 - p(1-s)\bigr)^{T}.
\]
\end{lemma}

\begin{proof}
Induction on $T$.  For $T=0$ both sides are $1$.  For the step, condition on
the first $T$ coordinates: with $G' = \sum_{t<T} X_t$,
\[
  \E\bigl[s^{G'+X_T}\bigr]
  = \E\Bigl[\, s^{G'} \cdot \E\bigl[s^{X_T} \mid \omega_0,\dots,\omega_{T-1}\bigr] \Bigr],
\]
and pointwise
$\E[s^{X_T} \mid \cdot\,] = 1 - q(1-s) \le 1 - p(1-s)$
where $q = \Pr[X_T = 1 \mid \cdot\,] \ge p$ and $1-s \ge 0$.  Since $s^{G'}
\ge 0$, the induction hypothesis applies to the outer expectation.
\end{proof}

\begin{corollary}[Lower tail]\label{cor:tail}
In the setting of Lemma~\ref{lem:adaptive}, for every $L \in \N$ and
$s \in (0,1]$,
\[
  \Pr[G < L] \;\le\; s^{-L}\,\bigl(1-p(1-s)\bigr)^{T}.
\]
With $p = \tfrac18$ and $s = \tfrac12$:
$\;\Pr[G < L] \le 2^{L}\,\bigl(\tfrac{15}{16}\bigr)^{T}$.
\end{corollary}

\begin{proof}
On $\{G < L\}$ we have $s^G \ge s^L$ (as $s \le 1$), so
$s^L \Pr[G < L] \le \E[s^G]$; apply Lemma~\ref{lem:adaptive} and Markov's
reasoning directly (no measure theory needed: these are finite sums).
\end{proof}

\begin{remark}
Lemma~\ref{lem:adaptive} replaces, in one induction, both the stochastic
domination of adaptive trials by a binomial and the Chernoff bound for the
binomial lower tail.  This is the main simplification that makes the growth
phase formalization-friendly.
\end{remark}

\section{Growth phase}\label{sec:growth}

\begin{definition}[Good rounds]\label{def:good}
For $t < T$ let
\[
  X_t \;=\; \mathbf{1}\Bigl[\, m_{t+1} \ge \tfrac98\, m_t
  \;\;\text{or}\;\; m_t > \tfrac n2 \,\Bigr] \in \{0,1\}.
\]
$X_t$ is a function of $(\omega_0,\dots,\omega_t)$.
\end{definition}

\begin{lemma}\label{lem:goodprob}
For every $t$ and every prefix $(\omega_0,\dots,\omega_{t-1})$,
$\Pr[X_t = 1 \mid \omega_0,\dots,\omega_{t-1}] \ge \tfrac18$.
\end{lemma}

\begin{proof}
The prefix determines $I_t$.  If $m_t > n/2$, then $X_t = 1$ with conditional
probability $1$.  Otherwise $1 \le m_t \le n/2$ (note $m_t \ge m_0 = 1$ by
Lemma~\ref{lem:mono}), and since $\omega_t$ is uniform on $R_n$ and
independent of the prefix, Lemma~\ref{lem:goodround} applies.
\end{proof}

\begin{lemma}[Enough good rounds force half-saturation]\label{lem:force}
Let $L \in \N$ satisfy $(9/8)^L > n/2$.  If $\omega$ is such that
$\sum_{t < T_1} X_t(\omega) \ge L$, then $m_{T_1}(\omega) > n/2$.
\end{lemma}

\begin{proof}
Suppose $m_{T_1} \le n/2$.  By monotonicity (Lemma~\ref{lem:mono}),
$m_t \le n/2$ for all $t \le T_1$.  Hence every good round $t < T_1$ satisfies
$m_{t+1} \ge \tfrac98 m_t$, and every round satisfies $m_{t+1} \ge m_t$.  By
induction on the number of good rounds, $m_{T_1} \ge (9/8)^{L} \, m_0 =
(9/8)^L > n/2$, a contradiction.
\end{proof}

\begin{lemma}[Phase 1]\label{lem:phase1}
With $T_1 = \ceil{117 \log n} + 23$,
\[
  \Pr\bigl[\, m_{T_1} \le \tfrac n2 \,\bigr] \;\le\; \frac1n .
\]
\end{lemma}

\begin{proof}
Set $L := \ceil{9 \log n} + 1$; since $\log\tfrac98 \ge \tfrac19$
(Lemma~\ref{lem:logbounds}), we have
$\log\bigl((9/8)^L\bigr) = L\log\tfrac98 \ge 9\log n \cdot \tfrac19 = \log n$,
hence $(9/8)^L \ge n > n/2$.  By Lemma~\ref{lem:force}
and Corollary~\ref{cor:tail} (with $p = 1/8$, $s = 1/2$, justified by
Lemma~\ref{lem:goodprob}),
\[
  \Pr\bigl[m_{T_1} \le \tfrac n2\bigr]
  \;\le\; \Pr\Bigl[\textstyle\sum_{t<T_1} X_t < L\Bigr]
  \;\le\; 2^{L}\,\bigl(\tfrac{15}{16}\bigr)^{T_1}.
\]
It remains to check $2^L (15/16)^{T_1} \le 1/n$, equivalently
\[
  T_1 \log\tfrac{16}{15} \;\ge\; L \log 2 + \log n .
\]
Using $\log\frac{16}{15} \ge \tfrac1{16}$, $\log 2 < 0.6932$ and
$L \le 9\log n + 2$:
\[
  L\log 2 + \log n \;\le\; (9 \log n + 2)\cdot 0.6932 + \log n
  \;\le\; 7.239 \log n + 1.387 ,
\]
while
\[
  T_1 \log\tfrac{16}{15} \;\ge\; \frac{117\log n + 23}{16}
  \;\ge\; 7.312 \log n + 1.437 . \qedhere
\]
\end{proof}

\begin{remark}[On constants]\label{rem:constants}
The clean statement to formalize is: \emph{for all $L, T_1$ with
$(9/8)^L > n/2$ and $2^L (15/16)^{T_1} \le 1/n$, we have
$\Pr[m_{T_1} \le n/2] \le 1/n$} --- the numerics above merely exhibit
admissible values $L = \ceil{9\log n} + 1$, $T_1 = \ceil{117\log n} + 23$;
any larger $T_1$ also works.
\end{remark}

\section{Saturation phase}\label{sec:saturation}

\begin{lemma}[Phase 2]\label{lem:phase2}
Let $A = \{m_{T_1} > n/2\} \subseteq \Omega$ (an event determined by the
first $T_1$ coordinates).  Then for every $s \ge 0$,
\[
  \E\bigl[\, U_{T_1 + s}\, \mathbf{1}_A \,\bigr]
  \;\le\; \Bigl(\frac23\Bigr)^{s}\, n .
\]
\end{lemma}

\begin{proof}
Induction on $s$.  For $s = 0$: $U_{T_1} \le n$ pointwise.  For the step,
condition on the first $T_1 + s$ coordinates: on $A$ we have $m_{T_1+s} \ge
m_{T_1} > n/2$ by monotonicity, so Lemma~\ref{lem:contract} applies to the
(uniform, independent) round $\omega_{T_1+s}$ and gives, pointwise on each
prefix in $A$,
\[
  \E\bigl[U_{T_1+s+1} \mid \omega_0,\dots,\omega_{T_1+s-1}\bigr]
  \;\le\; \tfrac23\, U_{T_1+s} .
\]
Multiplying by $\mathbf{1}_A$ (a function of the prefix), taking outer
expectations, and applying the induction hypothesis finishes the proof.
\end{proof}

\begin{corollary}\label{cor:phase2}
With $T_2 = \ceil{6 \log n}$:
$\;\Pr\bigl[\, U_T \ge 1 \text{ and } A \,\bigr] \le n\,(2/3)^{T_2} \le 1/n$.
\end{corollary}

\begin{proof}
$U_T$ is $\N$-valued, so by Markov (Lemma~\ref{lem:markov}) applied to
$U_T\mathbf{1}_A$ and Lemma~\ref{lem:phase2},
$\Pr[U_T \mathbf{1}_A \ge 1] \le (2/3)^{T_2} n$.
And $(2/3)^{T_2} n \le 1/n$ iff $T_2 \log\frac32 \ge 2\log n$, which holds
since $\log\frac32 \ge \tfrac13$ (Lemma~\ref{lem:logbounds}) and
$T_2 \ge 6\log n$.
\end{proof}

\section{Proof of Theorem~\ref{thm:main}}

With $T_1 = \ceil{117\log n} + 23$, $T_2 = \ceil{6\log n}$, $T = T_1+T_2$:
\[
  \Pr[I_T \ne V] \;=\; \Pr[U_T \ge 1]
  \;\le\; \Pr[\,\overline{A}\,] + \Pr[\,U_T \ge 1 \text{ and } A\,]
  \;\le\; \frac1n + \frac1n
\]
by Lemma~\ref{lem:phase1} and Corollary~\ref{cor:phase2}. \qed

\section{The Lean formalization}\label{sec:lean}

The argument above is fully formalized in the accompanying repository
(library \texttt{RumorSpread}, namespace \texttt{RumorPush}); the build
contains no \texttt{sorry} and the main theorem depends only on the three
standard axioms (\texttt{propext}, \texttt{Classical.choice},
\texttt{Quot.sound}).  Design notes:

\begin{itemize}[leftmargin=2em]
\item \textbf{Probability space.} No measure theory, no PMF, no ENNReal.
  \texttt{avg} is a finite sum divided by a cardinality, and
  \texttt{expList $\alpha$ T F} — the expectation of a trajectory functional
  $F$ over $T$ i.i.d.\ uniform draws — is defined \emph{by recursion on $T$}
  (average out the first round, recurse).  ``Conditioning on the first
  round'' is thereby a definitional unfolding, and Markov / union bounds are
  pointwise inequalities pushed through monotonicity of \texttt{expList}.
  The bridge theorem \texttt{expList\_eq\_avg\_ofFn}
  (\texttt{Equivalence.lean}) proves this equals the uniform average over
  the product space \texttt{Fin T $\to$ R n} of all round sequences, so the
  model is the textbook one.
\item \textbf{Process} (\texttt{Model.lean}): a round configuration is
  \texttt{R n := $\forall$ v : Fin n, \{u // u $\ne$ v\}}; one round is
  \texttt{step I r := I $\cup$ I.image (fun v => (r v).1)}; trajectories are
  lists of rounds consumed by \texttt{run}.  Lemma~\ref{lem:force} is the
  deterministic \texttt{pow\_goodCount\_mul\_card\_le\_card\_run}.
\item \textbf{Lemma~\ref{lem:contact}} (\texttt{avg\_not\_contacted} in
  \texttt{OneRound.lean}) is the only place where a probability is
  \emph{computed} rather than estimated: a counting argument via an
  explicit equivalence to a product of subtypes and
  \texttt{Fintype.card\_pi}.
\item \textbf{Concentration}: Lemma~\ref{lem:adaptive} becomes the
  ten-line induction \texttt{expList\_half\_pow\_goodCount}
  (\texttt{Growth.lean}); no Chernoff, no martingales, no couplings.
\item \textbf{Analytic toolkit} (\texttt{Bounds.lean}):
  \texttt{one\_add\_mul\_le\_pow} (Bernoulli) and short inductions for
  Lemmas~\ref{lem:bonferroni} and~\ref{lem:upper}.  \texttt{Real.log}
  enters only through the \emph{definition} of $L, T_1, T_2$ and the
  numeric checks of Section~\ref{sec:growth}, via
  Lemma~\ref{lem:logbounds} and Mathlib's
  \texttt{Real.log\_two\_lt\_d9}; \texttt{Real.exp} appears only inside the
  proof of \texttt{numeric\_A}.
\item \textbf{File structure:}
  \texttt{Bounds.lean} (Section~\ref{sec:ineq}),
  \texttt{Prob.lean} (\texttt{avg}/\texttt{expList}),
  \texttt{Model.lean} (Section~\ref{sec:model}),
  \texttt{OneRound.lean} (Section~\ref{sec:oneround}),
  \texttt{Growth.lean} (Sections~\ref{sec:adaptive}--\ref{sec:growth}),
  \texttt{Saturation.lean} (Section~\ref{sec:saturation}),
  \texttt{Main.lean} (main theorem
  \texttt{push\_informs\_all\_whp}),
  \texttt{Equivalence.lean} (faithfulness).
\end{itemize}

\begin{thebibliography}{9}
\bibitem{FG85} A.~Frieze, G.~Grimmett.
\emph{The shortest-path problem for graphs with random arc-lengths}.
Discrete Applied Mathematics 10 (1985) 57--77.
\bibitem{Pittel87} B.~Pittel.
\emph{On spreading a rumor}. SIAM J.\ Appl.\ Math.\ 47 (1987) 213--223.
\bibitem{KSSV00} R.~Karp, C.~Schindelhauer, S.~Shenker, B.~V\"ocking.
\emph{Randomized rumor spreading}. FOCS 2000.
\end{thebibliography}

\end{document}
